% bank: papit-token-pruning
% PAPIT -- the general GenAI pick: attention distillation, accuracy and scoring cost, and
% the alignment diagnostic. Facts and constraints: content/bank/projects/papit-token-pruning.tex.
% The title shows no "PAPIT:" prefix and the tech list says "LLaVA", both by request.
  \resumeProjectEntry
    {\textbf{\textcolor{RoyalBlue}{\href{https://github.com/yarikama/PAPIT}{Prompt-Aware Image Token Pruning for Multimodal Inference}}} $|$ \emph{PyTorch, LLaVA, CLIP ViT, Knowledge Distillation}}{Feb. 2026 -- May 2026}
    \resumeItemListStart
      \resumeItem{\textbf{Attention Distillation:} Taught a \textbf{1.3M-parameter} scorer to imitate where LLaVA-1.5-7B actually looks, dropping the patches it ignores before the projector: \textbf{10$\times$} less attention compute for the frozen LLM, and \textbf{150$\times$} cheaper than a CLIP-GradCAM pass}
      \resumeItem{\textbf{Accuracy \& Latency:} Kept a quarter of the image tokens and still answered TextVQA about as well as the full model, \textbf{12.3} points clear of random pruning, while deciding what to keep in \textbf{0.43 ms} an image}
      \resumeItem{\textbf{Empirical Diagnostic:} Showed why the training-free shortcut fails: across three VQA benchmarks, what CLIP finds salient barely tracks what LLaVA attends to, which is what made learning the scorer worth the trouble}
    \resumeItemListEnd%
